\chapter{The SDL\_RENDERER Backend}
\label{ch:sdl-renderer}

\texttt{SDL\_RENDERER} is CNA's simplest backend, its most exhaustively audited, and the one
whose scope is narrowest by explicit design: a 2D-only rendering path built directly on
SDL3's own 2D texture-blit renderer, with no 3D pipeline, no programmable shader stage, no
depth/stencil buffer, and no multisample antialiasing at all.

\section{Scope, by design}

Every 3D-facing entry point on this backend --- \texttt{CreateVertexBuffer},
\texttt{DrawColoredPrimitives}, and the rest of the 3D draw dispatch from
Chapter~\ref{ch:backend-architecture} --- throws immediately, with a consistent
\texttt{"SDL\_Renderer does not support 3D"} message. This is not treated as a bug: a
systematic full-suite run confirms exactly thirteen known, expected pre-existing test
failures, every one of them this exact throw, matching the backend's documented scope
precisely. A custom \cnaclass{Effect} passed to \texttt{SpriteBatch::Begin(effect)} also
throws, for the same underlying reason --- there is no shader stage to run it on.

\subsection{A worked example: the one backend that answers every capability honestly}

Chapter~\ref{ch:backend-architecture}'s \cnaclass{GraphicsCapability} mechanism defaults to
answering \texttt{true} for everything unless a backend overrides it --- and reading
\texttt{SdlGraphicsBackend.hpp} directly shows this backend takes the blanket, honest opposite
stance rather than enumerating exceptions one at a time:

\begin{lstlisting}
bool SupportsCapability(GraphicsCapability /*capability*/) const override
{
    // 2D-only by design: none of the capabilities GraphicsCapability currently
    // enumerates are supported on this backend.
    return false;
}
\end{lstlisting}

Every single \cnaclass{GraphicsCapability} value --- not just \texttt{ThreeD}, but
\texttt{DepthStencilBuffer}, \texttt{Multi\allowbreak{}Sample\allowbreak{}AntiAliasing},
\texttt{OcclusionQuery}, all
eight --- returns \texttt{false} on \texttt{SDL\_RENDERER}, unconditionally. This makes it the
one backend in the whole project where \texttt{SupportsCapability} genuinely never needs a
per-feature audit the way Chapter~\ref{ch:backend-architecture}'s own
\cnaclass{MultipleRenderTargets} finding showed was necessary elsewhere --- a real, useful
asymmetry to know before trusting any other backend's answer the same way:

\begin{lstlisting}
if (device.SupportsCapability(GraphicsCapability::ThreeD)) {
    DrawModel(model, world, view, projection);
} else {
    // Reached on SDL_RENDERER (and CANVAS) -- checking first here is strictly better
    // than letting CreateVertexBuffer() or DrawColoredPrimitives() throw
    // "SDL_Renderer does not support 3D" and catching it after the fact.
    DrawFlatPlaceholderSprite();
}
\end{lstlisting}

\section{The Phase 70 audit: fifteen real bugs, two still open}

The single most thorough backend audit in the project (Tasks~666--731, 65 tasks total)
found and fixed fifteen real bugs, spanning almost every area \cnaclass{SpriteBatch}
touches:

\begin{itemize}
  \item \textbf{SpriteBatch (2 bugs).} \texttt{SDL\_RenderTextureRotated}'s rotation-pivot
        convention did not match XNA's own --- fixed by offsetting the destination
        rectangle. The \texttt{transformMatrix} argument to \texttt{Begin()} was silently
        ignored in its entirety --- fixed via a new path using
        \texttt{SDL\_RenderTextureAffine}, used only for a genuinely non-identity transform
        (Chapter~\ref{ch:spritebatch} covers both in full).
  \item \textbf{SpriteFont (1 bug, shared across every backend).} The
        \cnaclass{SpriteEffects} flip bug described in Chapter~\ref{ch:spritebatch} ---
        fixed once, in the shared \texttt{SpriteBatch.cpp}, benefiting every backend at
        once rather than being an SDL\_Renderer-specific fix.
  \item \textbf{BlendState (2 bugs).} Only \texttt{Opaque} and \texttt{Additive} were
        specially handled; every other preset mapped incorrectly. Separately,
        \texttt{SpriteBatch::Begin()} was unconditionally clobbering whatever blend mode had
        already been set, every single call.
  \item \textbf{SamplerState (1 bug).} Four of nine \cnaclass{TextureFilter} values were
        silently downgraded to nearest-neighbor sampling rather than mapping to their real
        \texttt{SDL\_ScaleMode} equivalent.
  \item \textbf{RenderTarget2D (4 bugs).} A shared depth-clear code path crashed outright
        when asked to clear the depth buffer of a render target with
        \texttt{DepthFormat::None} (i.e., no depth buffer at all, which is every render
        target on this 2D-only backend). A second bug performed an unchecked, unsafe
        downcast between two sibling backend classes when sampling a render target as an
        ordinary texture --- a genuine memory-safety hazard, not merely a logic error. The
        requested \texttt{DepthFormat} itself is honestly acknowledged as emulated: it is
        silently accepted and echoed back, but no functional depth test exists behind it at
        all on this backend. A new \texttt{HasRealDepthBuffer()} check was added specifically
        so calling code can ask, rather than assume.
  \item \textbf{Viewport (1 bug).} \texttt{PresentInterval::Two} silently collapsed to
        \texttt{PresentInterval::One} rather than being honored or rejected.
  \item \textbf{Disposed-resource guards (3 bugs).} A missing disposed-texture check inside
        the internal sprite-push path; a missing equivalent check in
        \texttt{SetRenderTargets}; and a dangling backend pointer left behind after
        \texttt{RenderTarget2D::Dispose()}.
  \item \textbf{OcclusionQuery (1 bug).} Construction previously succeeded silently even with
        a null backend --- now throws, consistent with the 2D-only-backend pattern of
        failing loudly at construction rather than producing a query object that can never
        do anything.
\end{itemize}

\section{Two blocked decisions, and why they are genuinely blocked}

Two gaps remain open specifically because closing them requires a project-owner
architecture decision, not more engineering effort:

\begin{description}
  \item[\texttt{TextureAddressMode::Wrap} / \texttt{Mirror} via \texttt{SpriteBatch}.]
        \texttt{SDL\_RenderTexture}'s own \texttt{srcrect} handling has exactly one fixed,
        clamp-like behavior at the edges. SDL3 does expose
        \texttt{SDL\_SetRenderTextureAddressMode}, but it only affects
        \texttt{SDL\_RenderGeometry} --- an entirely different draw call than the one this
        backend's \texttt{Draw()} actually issues. Three concrete options were identified and
        none chosen: throw unconditionally whenever \texttt{Wrap} / \texttt{Mirror} is
        requested, rewrite \texttt{Draw()} onto \texttt{SDL\_RenderGeometry} project-wide, or
        a hybrid that only throws when the source rectangle actually exceeds the texture's
        bounds (the common case where it does not would keep working, coincidentally, under
        the existing clamp-like behavior).
  \item[\texttt{Texture3D} / \texttt{TextureCube} construction.] These currently succeed
        silently with a null backend rather than throwing. Prototyping a fix that makes
        construction throw instead was found to have a 94-test blast radius (33 tests in one
        file, 42 in another, 19 in a third) --- large enough that the change was left
        unresolved rather than made unilaterally.
\end{description}

\subsection{A worked example: exactly what ``clamp-like, by coincidence'' means in pixels}

The blocked \texttt{TextureAddressMode::Wrap} / \texttt{Mirror} decision above is easy to read
as an abstract API gap; a concrete pixel case makes it precise instead. Take a $2\times1$
texture, texel~0 red and texel~1 blue, and draw it with a source rectangle \emph{twice} the
texture's own width --- source texel range $[0,4]$ against a 2-texel-wide texture, the classic
XNA scrolling-background tiling technique, which FNA's real \cnaclass{SpriteBatch} never clamps
\texttt{sourceRectangle} to the texture's own bounds to prevent:

\begin{lstlisting}
Texture2D redBlue(getGraphicsDeviceProperty(), 2, 1);
const Color pixels[2] = {Color::Red, Color::Blue};
redBlue.SetData(pixels, 2);

Rectangle oversizedSource(0, 0, 4, 1); // spans texel [0,4] on a 2-texel-wide texture
spriteBatch_->Begin(SpriteSortMode::Deferred, BlendState::Opaque, &pointClamp, nullptr, nullptr);
spriteBatch_->Draw(redBlue, destRect, oversizedSource, Color::White);
spriteBatch_->End();
\end{lstlisting}

Reading back the destination pixel corresponding to source position $1.25$ (three-quarters of
the way into the oversized range, past the real 2-texel-wide texture entirely) is where the
``clamp-like, by coincidence'' finding becomes a concrete, checkable number rather than a
qualitative claim: with \texttt{PointClamp} requested, the correct XNA answer is
\textbf{blue} --- clamped to the texture's last real texel --- and \texttt{SDL\_RenderTexture}'s
one fixed edge behavior genuinely produces exactly that, with zero production code written to
make it happen. Requesting \texttt{PointWrap} at the identical source position is where the
still-open gap becomes equally concrete: the XNA-correct answer for a wrapped read at position
$1.25$ is \textbf{red} (wrapping back around to texel~0), but this backend's \texttt{Draw()}
produces \textbf{blue} again --- the same fixed clamp-like behavior applies regardless of which
\cnaclass{TextureAddressMode} was actually requested, since (per the finding above)
\texttt{SDL\_RenderTexture} has no address-mode parameter reachable from this specific draw
path at all. A game relying on horizontally-scrolling tiled backgrounds --- a common enough
technique that this is not a hypothetical edge case --- gets the wrong pixel on this backend
today, silently, with no exception and no visual artifact obviously identifiable as ``wrong''
unless you already know to look for texel~0 failing to reappear at the wrap boundary.

\section{What is honestly emulated versus what is real}

Beyond the \texttt{DepthFormat} case above, \texttt{TextureAddressMode::Clamp} is flagged as
working ``by coincidence'' --- it matches \texttt{SDL\_RenderTexture}'s own fixed
out-of-bounds behavior, not because any real sampler-state value is being honored.
\texttt{MultiSampleCount} is accepted and ignored, always reporting zero. Fullscreen
toggling round-trips correctly at the API level, but the audit is explicit that genuine
OS-level fullscreen state cannot actually be verified inside this project's own headless
\texttt{Xvfb} test environment. \texttt{Clear} honoring \texttt{ClearOptions::Stencil} is a
cross-backend no-op here too (Task~871, still open, matching Chapter~\ref{ch:state-objects}'s
finding for the hardware backends). Five real, multi-frame compatibility samples --- a 2D
demo, a bouncing-sprite physics scene, a keyboard-driven sprite, a two-glyph
\cnaclass{SpriteFont} test, and an animated spritesheet --- are registered as genuine
\texttt{ctest} entries rather than left as manual verification, giving this backend
end-to-end coverage beyond its unit-level pixel tests.

\section{Reading the feature matrix for this backend}

\texttt{docs/graphics-backend-feature-matrix.md}'s own snapshot for
\texttt{SDL\_RENDERER} shows a 67-of-67 integration-test pass rate; the custom-\cnaclass{Effect}
throw-by-design is the matrix's own explicit gap marker, alongside SDL\_Renderer's two named
BLOCKED rows (\texttt{Wrap} / \texttt{Mirror}, and \texttt{Texture3D} / \texttt{TextureCube}
construction) --- the same two items covered above, and, notably, two of only four
project-wide entries anywhere in the whole feature matrix that carry the ``blocked, needs a
project-owner decision'' status rather than either a solid pass or an openly acknowledged,
still-being-worked gap.
